\chapter{About This Manual}

This manual describes how to operate the \productname\ \productdescriptor: its
controls, its menus, its bonding modes, and the parameters that determine bond
quality.

\section*{Who it is for}

The operator at the machine. It assumes you understand what a wire bond is and
why one might be good or bad, but not that you have used this particular
machine before.

It does not cover installation, mechanical service, electrical repair, or
firmware development. Those belong in a service manual.

\section*{How it is organised}

\begin{description}[leftmargin=!,labelwidth=42mm,style=nextline,font=\normalfont\bfseries]
\item[Part I --- Getting to Know the Machine]
  The axes, the controls, the display, and what happens at power-up. Read this
  first.
\item[Part II --- Operating]
  Navigating the menus, managing configurations, the four bonding modes, and
  running a bond.
\item[Part III --- Setup and Tuning]
  The three calibrations you perform from the panel, and how to turn a bad bond
  into a good one.
\item[Part IV --- Reference]
  Every parameter, every message, the specifications, and a one-page quick
  reference to print and keep at the machine.
\end{description}

\section*{Conventions}

\begin{description}[leftmargin=!,labelwidth=36mm,style=nextline]
\item[\keycap{SAVE}] A physical button on the panel.
\item[\menuitem{Loop Height}] A label or value as it appears on the display.
\item[\msg{OPERATION FAILED}] A message the machine shows you.
\end{description}

Three kinds of aside appear throughout:

\begin{wbnote}
Something worth knowing that is easy to miss.
\end{wbnote}

\begin{wbwarning}
Something that can damage the machine, the workpiece, or you.
\end{wbwarning}

\begin{wbtodo}
A section not yet written. None of these should remain in a released manual.
\end{wbtodo}
